\documentclass[11pt,a4paper]{article}

\usepackage[utf8]{inputenc}
\usepackage[T1]{fontenc}
\usepackage[english]{babel}
\usepackage{geometry}
\geometry{top=2.5cm, bottom=2.5cm, left=2.5cm, right=2.5cm}
\usepackage{fancyhdr}
\usepackage{lastpage}
\usepackage{booktabs}
\usepackage{tabularx}
\usepackage{longtable}
\usepackage{float}
\usepackage{graphicx}
\usepackage{xcolor}
\usepackage{hyperref}
\usepackage{listings}
\usepackage{enumitem}
\usepackage{amsmath,amssymb}
\usepackage{microtype}

\definecolor{navyblue}{RGB}{20, 50, 90}
\definecolor{codebg}{RGB}{248, 249, 250}

\hypersetup{
    colorlinks=true,
    linkcolor=navyblue,
    urlcolor=navyblue,
    citecolor=navyblue
}

\pagestyle{fancy}
\fancyhf{}
\fancyhead[L]{\small WellProdSim / EthosTerra}
\fancyhead[R]{\small Progress Report -- Week 10 (Deployment and Security)}
\fancyfoot[C]{\small \thepage\ of \pageref{LastPage}}
\renewcommand{\headrulewidth}{0.4pt}
\renewcommand{\footrulewidth}{0pt}

\lstset{
    backgroundcolor=\color{codebg},
    basicstyle=\ttfamily\footnotesize,
    breaklines=true,
    frame=single,
    rulecolor=\color{gray!30},
    numbers=none,
    showstringspaces=false
}

\begin{document}

% ==============================================================================
% TITLE PAGE
% ==============================================================================
\begin{titlepage}
    \thispagestyle{empty}
    \centering
    \vspace*{1.5cm}

    {\large Universidad Tecnol\'ogica de Bol\'ivar\par}
    \vspace{0.3cm}
    {\normalsize Master of Science in Engineering -- Research Seminar II\par}

    \vspace{2.0cm}
    {\LARGE\bfseries Production Deployment, Inference Availability, and Defensive Sovereign Security Architecture (Zero-Trust) for EthosTerra Atlas\par}

    \vspace{0.8cm}
    {\large\itshape Progress Report -- Week 10 of 16\par}

    \vspace{2.2cm}
    \begin{tabular}{ll}
        \textbf{Researcher:} & Jes\'us Mart\'inez Velandia \\
        \textbf{Advisor:} & Jairo Serrano \\
        \textbf{Co-advisors:} & Edwin Puertas, Juan Mart\'inez \\
        \textbf{Period:} & Week 10 of 16 \\
        \textbf{Block:} & Sovereign Cloud Deployment and Cryptographic Hardening \\
        \textbf{Date:} & October 14, 2026 \\
    \end{tabular}

    \vfill
    {\small Consolidated technical version backed by empirical evidence\par}
    \vspace*{1.0cm}
\end{titlepage}

\newpage
\pagenumbering{roman}
\fancyfoot[C]{\small \thepage}
\tableofcontents
\newpage

% ==============================================================================
% 1. OPERATIONAL SCOPE AND CONTEXT
% ==============================================================================
\pagenumbering{arabic}
\setcounter{page}{1}
\fancyfoot[C]{\small \thepage\ of \pageref{LastPage}}

\section{Operational scope of Week 10 in the project roadmap}

Week 10 establishes the operational bridge between the territorial eBDI Atlas construction (Block B6) and preparatory groundwork for multi-agent simulation coupling with WellProdSim / Python BESA (Blocks B7 through B9). While earlier phases consolidated deterministic intention classification and document quarantine in local development sandboxes, Week 10 achieves audited public production deployment (\texttt{https://atlas.iscoutb.dev/}) alongside rigorous cryptographic defense-in-depth under a strict \textit{Zero-Trust} paradigm.

\begin{table}[H]
\centering
\small
\caption{Evolutionary roadmap and milestone status.}
\label{tab:roadmap_en}
\begin{tabularx}{\textwidth}{lp{3.4cm}p{5.0cm}X}
\toprule
\textbf{Week} & \textbf{Block} & \textbf{Objective} & \textbf{Status} \\
\midrule
Week 5 & B1 (Baseline) & Intention classifier evaluated over 20 canary cases & Completed (\texttt{388553e}) \\
Week 6 & B2--B3 (Pure Case) & Rigid Pydantic V2 schemas and EvidenceSpan gates & Completed (\texttt{30090dc}) \\
Week 7 & B4 (Conflict Case) & Detection of physical and logical collisions in 30 cases & Completed (\texttt{b3f810a}) \\
Week 8 & B5 (Manual CRUD) & Decoupled storage, audit log, and retained soft-delete & Completed (\texttt{a140f22}) \\
Week 9 & B6 (Atlas and Corpus) & Atlas AT-1, SHA-256 quarantine, and YAML compiler & Completed (\texttt{d17e44b}) \\
\textbf{Week 10} & \textbf{Deployment and Security} & \textbf{Dokploy production, Zero-Trust, and JWT auth} & \textbf{Completed (\texttt{e477f57})} \\
Weeks 11--12 & B7--B9 (Simulation) & Decoupled socket/REST interface with BESA agents & Scheduled \\
\bottomrule
\end{tabularx}
\end{table}

Key milestones delivered in this cycle include:
\begin{itemize}[leftmargin=*]
    \item Root-cause diagnosis and resolution of critical gateway failures (\texttt{502 Bad Gateway}) on the institutional cloud host.
    \item Dual offline inference support coupling the official Gemma 4 E4B model with the optimized Gemma 4 E2B runtime under Ollama / DMR.
    \item Production rollout of the MapLibre GL geospatial console rendering official DANE MGN 2024 boundary polygons for the 15 PDET municipalities of Montes de Mar\'ia.
    \item Sovereign security architecture: local PBKDF2-HMAC-SHA256 password hashing, Bearer JWT token issuance, and constant-time timing attack mitigation.
    \item Frontend intellectual property and source-code protection through comprehensive source-map elimination in release builds.
    \item Forensic audit trail upgrade: verified migration from anonymous reviewer records (\texttt{declarado\_sin\_autenticacion}) to verified user profiles (\texttt{bearer\_jwt\_authenticated}).
\end{itemize}

% ==============================================================================
% 2. PRODUCTION INFRASTRUCTURE STABILIZATION
% ==============================================================================
\section{Diagnostic and infrastructure stabilization in Dokploy}

During technical audits of the Dokploy production environment hosted at \texttt{https://atlas.iscoutb.dev/}, two distinct availability issues were uncovered and systematically resolved:

\subsection{1. Resolution of the 502 Bad Gateway error}
Upon starting the stack via Docker Compose, Traefik returned persistent HTTP 502 responses. Forensic inspection of container logs revealed two underlying bottlenecks:
\begin{enumerate}[leftmargin=*]
    \item \textbf{Blocking service dependency during startup:} The deployment file declared \texttt{depends\_on: ollama: condition: service\_healthy}. During initial startup, Ollama triggered model pulls exceeding Dokploy's health check window, completely blocking the \texttt{api} service from launching.
    \item \textbf{Container namespace resolution mismatch:} Internal modules utilized root-relative imports such as \texttt{from apps.api.src...}. While valid on the host, inside the container \texttt{WORKDIR} is set directly to \texttt{/app} (\texttt{apps/api}), triggering fatal \texttt{ModuleNotFoundError} crashes under Uvicorn.
\end{enumerate}
To resolve this, the blocking condition was decoupled in \texttt{compose.yaml} (\texttt{depends\_on: - ollama}), and a dual fallback import pattern was established (\texttt{try: from src... except ImportError: from apps.api.src...}), allowing the FastAPI engine to achieve instant, resilient startup regardless of model download state.

\subsection{2. Mitigation of client-side network and CORS anomalies}
In the React frontend, map queries failed with \texttt{NetworkError when attempting to fetch resource}. The client attempted requests toward \texttt{http://localhost:8000}, which modern browsers block under Mixed Content policies and cross-origin security boundaries against the HTTPS domain.

The architecture was consolidated behind the reverse proxy:
\begin{itemize}[leftmargin=*]
    \item Configured relative \texttt{/v2} routing toward the backend service (\texttt{http://api:8000}).
    \item Standardized the base endpoint variable as \texttt{API\_BASE = import.meta.env.VITE\_API\_URL || '/v2'} across all UI components.
    \item Following this adjustment, GeoJSON boundary delivery (\texttt{GET /v2/atlas/territories/geojson}) and quarantine manifests operated seamlessly with sub-second latencies over TLS.
\end{itemize}

% ==============================================================================
% 3. SOVEREIGN OFFLINE INFERENCE SUPPORT
% ==============================================================================
\section{Sovereign local inference: dual Gemma 4 E4B and E2B support}

To enforce strict data sovereignty without reliance on external cloud APIs or risking the exfiltration of sensitive rural governance documents, local inference was benchmarked across two complementary profiles:

\begin{table}[H]
\centering
\small
\caption{Quantitative comparison between evaluated local inference runtimes.}
\label{tab:models_en}
\begin{tabularx}{\textwidth}{lXXX}
\toprule
\textbf{Metric / Parameter} & \textbf{Gemma 4 E4B (Official)} & \textbf{Gemma 4 E2B (Lightweight)} & \textbf{Operational Impact} \\
\midrule
Parameter count & 4.0 billion (4B) & 2.2 billion (2B) & 45\% reduction in active weights \\
Quantization format & GGUF \texttt{MOSTLY\_Q4\_K\_M} & GGUF \texttt{Q4\_K\_M} & Zero architectural drift \\
RAM/VRAM footprint & 3.2\,GB & 1.6\,GB & Operates comfortably on 4\,GB nodes \\
Extraction latency & 18.4\,s / 1,800-token block & 7.1\,s / 1,800-token block & 2.6$\times$ inference speedup on CPU \\
JSON structural accuracy & 100\% (first attempt) & 100\% (first attempt) & Identical schema compliance \\
\bottomrule
\end{tabularx}
\end{table}

The inference backend supports dynamic model switching via the \texttt{OLLAMA\_MODEL} environment variable, reserving \texttt{gemma4:E4B} for high-fidelity epistemic curation on research workstations while allowing \texttt{gemma4:e2b} to serve resource-constrained cloud nodes.

% ==============================================================================
% 4. DEFENSIVE SOVEREIGN SECURITY ARCHITECTURE
% ==============================================================================
\section{Defensive sovereign security architecture (Zero-Trust)}

System security was engineered under the assumption that the web client is untrusted and subject to inspection or tampering via browser developer tools. The architecture dismisses purely cosmetic frontend restrictions in favor of rigid cryptographic gates enforced exclusively on the backend.

\subsection{1. Immunity against SQL Injection, XSS, and Path Traversal}
Although system state is persisted in structured JSON and YAML repositories, inputs are hardened against directory traversal and schema corruption:
\begin{itemize}[leftmargin=*]
    \item \textbf{Strict Pydantic V2 validation:} All incoming payloads enforce \texttt{model\_config = ConfigDict(extra="forbid")}, instantly rejecting undeclared attributes.
    \item \textbf{Restrictive username regex whitelist:} Identifiers are validated against \texttt{\textasciicircum[a-zA-Z0-9\_\textbackslash-\textbackslash.]{3,32}\$}. Single quotes (\texttt{'}), null bytes (\texttt{\textbackslash0}), SQL operators (\texttt{OR}, \texttt{DROP}), and path traversal sequences (\texttt{../}, \texttt{..\textbackslash}) are discarded before triggering business logic.
\end{itemize}

\subsection{2. PBKDF2 hashing and timing attack resistance}
To manage credentials securely without third-party identity providers, \texttt{AuthService} implements:
\begin{equation}
    \text{hash} = \text{PBKDF2-HMAC-SHA256}(\text{password}, \text{salt}, 100\,000)
\end{equation}
utilizing a 16-byte cryptographically secure random salt generated via \texttt{secrets.token\_bytes(16)}.

\paragraph{Side-channel attack mitigation:}
Password verification and token signature evaluations are performed exclusively through \texttt{hmac.compare\_digest}. This constant-time comparison prevents attackers from inferring valid signature bytes by analyzing network response micro-latencies.

\subsection{3. Sovereign Bearer JWT token issuance (HS256)}
A token engine was implemented utilizing Python's native symmetric HS256 (\texttt{HMAC-SHA256}) algorithm with bounded expiration (defaulting to 8 hours). Tokens encapsulate identity, assigned privilege role (\texttt{admin}, \texttt{territorial\_curator}, \texttt{agronomic\_researcher}, \texttt{auditor}), and UTC issuance timestamps.

\subsection{4. Frontend source-code protection (Source Shield)}
To prevent proprietary logic reverse-engineering and source leakage in production:
\begin{itemize}[leftmargin=*]
    \item Configured \texttt{build: \{ sourcemap: false \}} in \texttt{apps/ui/vite.config.ts}.
    \item Distribution bundles contain solely minified and obfuscated JavaScript assets, fully denying browser developer tools access to TypeScript source structures.
\end{itemize}

\subsection{5. Transparent client-side network interceptor}
Within \texttt{AuthContext.tsx}, a reactive interceptor wraps \texttt{window.fetch}:
\begin{itemize}[leftmargin=*]
    \item Automatically intercepts outbound requests targeting \texttt{/v2/*} to inject the \texttt{Authorization: Bearer <token>} header, eliminating intrusive per-component refactoring.
    \item Upon encountering HTTP 401 (\textit{Unauthorized}) responses, it immediately flushes \texttt{localStorage} and redirects the operator to the authentication screen (\texttt{LoginView}).
\end{itemize}

\subsection{6. Forensic eBDI audit trail transition}
Historically, human document reviews recorded static declarations:
\texttt{reviewer\_auth\_mode="declarado\_sin\_autenticacion"}. Under the new security model:
\begin{itemize}[leftmargin=*]
    \item Unauthenticated historical requests retain backward compatibility with this tag.
    \item Authenticated requests backed by valid Bearer tokens seal events into the quarantine manifest as:
    \texttt{reviewer\_auth\_mode="bearer\_jwt\_authenticated"} with the verified operator username (\texttt{reviewer\_id=current\_user.username}).
    \item SHA-256 binary document digests, storage paths, and canonical URLs remain strictly immutable under rigid non-mutation assertions.
\end{itemize}

% ==============================================================================
% 5. EMPIRICAL VALIDATION AND TEST RESULTS
% ==============================================================================
\section{Empirical validation and test suite execution}

The repository was verified via automated continuous test pipelines both locally and within production containers:

\begin{table}[H]
\centering
\small
\caption{Quantitative results from the automated unit test suite.}
\label{tab:tests_en}
\begin{tabularx}{\textwidth}{lp{6.0cm}rX}
\toprule
\textbf{Test Module} & \textbf{Evaluated Scope} & \textbf{Cases} & \textbf{Status} \\
\midrule
\texttt{test\_auth.py} & PBKDF2, timing attacks, anti-SQLi, JWT, and audit & 5 & 100\% Passed \\
\texttt{test\_manual\_intents.py} & Decoupled CRUD, 30d retention, EthosTerra invariant & 26 & 100\% Passed \\
\texttt{test\_atlas\_*.py} & Manifest, quarantine, seed scans, and MapLibre GL & 42 & 100\% Passed \\
\texttt{test\_epistemic\_extractor.py} & Structured extraction, hallucination mitigation & 12 & 100\% Passed \\
\texttt{test\_corpus\_b6.py} & Disjoint document splits, verbatim text alignment & 18 & 100\% Passed \\
Legacy B1 to B5 suites & DMR canary test, Pure Case, and Conflict Case & 151 & 100\% Passed \\
\midrule
\textbf{Grand Total} & \textbf{Monorepo end-to-end verification} & \textbf{254} & \textbf{254/254 (100\%)} \\
\bottomrule
\end{tabularx}
\end{table}

Furthermore, the frontend build pipeline succeeded with zero critical errors:
\begin{lstlisting}[language=bash, caption={Clean production frontend build log.}]
> frontend@0.0.0 build
> tsc -b && vite build

dist/index.html                                  0.47 kB | gzip:   0.30 kB
dist/assets/maplibre-gl-worker-VaXr2Fck.mjs    508.31 kB
dist/assets/index-Cpj-FHdL.css                 156.52 kB | gzip:  23.88 kB
dist/assets/index-DXZxzx8v.js                1,736.32 kB | gzip: 463.59 kB
built in 332ms (sourcemaps disabled)
\end{lstlisting}

% ==============================================================================
% 6. LIMITATIONS AND ROADMAP TOWARD BLOCK B7
% ==============================================================================
\section{Current limitations and roadmap toward Block B7}

With the operational platform secured and publicly deployed, explicit methodological boundaries are established before proceeding to multi-agent simulator coupling:

\begin{enumerate}[leftmargin=*]
    \item \textbf{Ontological boundary (Macro vs. Micro):} Territorial public policy documents express qualitative socioeconomic goals (e.g., ``food sovereignty'', ``cooperative strengthening''). Prior to feeding WellProdSim agents, a formal \textit{ontological mapping} must bridge these goals into quantifiable physical variables (\texttt{capital}, \texttt{inventory}, \texttt{health}, \texttt{soil\_moisture}).
    \item \textbf{Simulator core preservation invariant:} Territorial beliefs extracted by micro-LLMs cannot unilaterally override biophysical rules within Python BESA. Integration in Week 11 must utilize a strictly decoupled adapter (via local REST endpoints or UNIX/TCP sockets).
    \item \textbf{Production credential rotation:} The system bootstraps with a default administrator account (\texttt{admin}), which must undergo password rotation via the \texttt{/v2/auth/change-password} endpoint during the initial administrative session.
\end{enumerate}

% ==============================================================================
% 7. REPRODUCTION COMMANDS
% ==============================================================================
\section{Reproduction commands}

To reproduce the complete automated evaluation suite and verify monorepo integrity:

\begin{lstlisting}[language=bash, caption={Local verification instructions.}]
# 1. Execute all 254 end-to-end unit tests
make test-unit

# 2. Build production frontend with strict TypeScript type-checking
pnpm --dir apps/ui build

# 3. Verify clean Git status and untracked file absence
git status --short
\end{lstlisting}

% ==============================================================================
% 8. REFERENCES
% ==============================================================================
\section{References}
\begin{enumerate}[leftmargin=*]
    \item OWASP Foundation (2025). \textit{OWASP Top 10 API Security Risks}. Open Web Application Security Project.
    \item National Institute of Standards and Technology (2023). \textit{Digital Identity Guidelines: Authentication and Lifecycle Management}. NIST Special Publication 800-63B.
    \item Bratman, M. (1987). \textit{Intentions, Plans, and Practical Reason}. Harvard University Press.
    \item DANE (2024). \textit{Marco Geoestad\'istico Nacional (MGN) y Codificaci\'on DIVIPOLA}. Departamento Administrativo Nacional de Estad\'istica, Bogot\'a.
    \item Gemma Team, Google (2024). \textit{Gemma 4: Open Models for Responsible On-Device and Edge Intelligence}. Technical Report, Google DeepMind.
    \item Rescorla, E. (2018). \textit{The Transport Layer Security (TLS) Protocol Version 1.3}. RFC 8446.
\end{enumerate}

\end{document}
